\documentclass[12pt, a4paper]{article}

% --- Margins & Layout ---
\usepackage[margin=1in, headheight=15pt]{geometry}

% --- Micro-Typography Adjustments ---
\usepackage{microtype} % Handles character protrusion and font expansion automatically
\usepackage[brazilian]{babel}

% --- Math & Symbols ---
\usepackage{amsmath, amssymb, amsthm}

% --- Built-in Header/Footer Control (No fancyhdr needed) ---
% Uses standard LaTeX page styles or lightweight native commands
\pagestyle{myheadings}
\markright{\scriptsize Pedro Barros \quad|\quad Medida e Integração \quad|\quad Sigma-álgebra e Medidas}

% --- Custom Math Environments ---
\newenvironment{solution}[1]{\par\vspace{1.2em}\noindent\textbf{Questão #1.}}{\par\vspace{0.5em}}

% --- Optional Math Shortcuts ---
\newcommand{\R}{\mathbb{R}}
\newcommand{\N}{\mathbb{N}}
\newcommand{\Z}{\mathbb{Z}}
\newcommand{\salg}{$\sigma$-álgebra }
\newcommand{\A}{\mathcal{A}}
\newcommand{\B}{\mathcal{B}}
\newcommand{\vp}{\varphi}

\begin{document}

% --- Title Block ---
\thispagestyle{empty} % Hides headers on the first page
\begin{center}
    {\Large \textbf{Introdução à Medida e Integração}} \\[0.4em]
    {\large Sigma-álgebra e Medidas} \\[0.4em]
    \textbf{Nome:} Pedro Barros \quad | \quad \textbf{Data:} \today
\end{center}

\hrule
\vspace{1.5em}

% --- Content ---
\begin{solution}{1}
    Seja $\lambda(E)=\mu(A\cap E)$, para todo $E\in\mathcal{A}$. Verifiquemos que $\lambda$ é uma medida.

    \textbf{Axioma 1:} Como $A\cap\varnothing=\varnothing$,
    \[
        \lambda(\varnothing)=\mu(A\cap\varnothing)=\mu(\varnothing)=0.
    \]

    \textbf{Axioma 2:} Para todo $E\in\mathcal{A}$, temos $A\cap E\in\mathcal{A}$. Logo,
    \[
        \lambda(E)=\mu(A\cap E)\geq 0.
    \]

    \textbf{Axioma 3:} Se $(E_n)_{n\geq1}\subseteq\mathcal{A}$ é uma família dois a dois disjunta, então os conjuntos $A\cap E_n$ também são dois a dois disjuntos e
    \[
        A\cap\biggl(\bigcup_{n=1}^{\infty}E_n\biggr)
        =\bigcup_{n=1}^{\infty}(A\cap E_n).
    \]
    Pela $\sigma$-aditividade de $\mu$,
    \[
        \lambda\biggl(\bigcup_{n=1}^{\infty}E_n\biggr)
        =\mu\biggl(A\cap\bigcup_{n=1}^{\infty}E_n\biggr)
        =\sum_{n=1}^{\infty}\mu(A\cap E_n)
        =\sum_{n=1}^{\infty}\lambda(E_n).
    \]
    Portanto, $\lambda$ é uma medida em $(X,\mathcal{A})$.
\end{solution}

\end{document}